\SourceLeaf{173}\noindent\hspace*{\parindent}Den Goldcours habe ich angenommen 1 Louisd’or \(=\) 9 fl. 38 kr. und \(=5\tfrac{1}{2}\) Thaler Courant. Theilweise sind übrigens die Course seitdem noch etwas, obwohl nur wenig, gewichen.

Das Resultat dieser Reductionen ist, dass dieselben verzinslichen Capitale, welche nach dem Nennwerthe zu 123372 Thaler angesetzt sind, nach dem zeitigen Börsencourswerthe nur 119044 Thl. 23 Ggr. erbringen, welcher Summe ich noch diejenigen 800 Thaler zusetze, für welche das vormals Müllersche Haus verkauft ist. So stellt sich unter Beifügung der andern Posten (baarer Geldvorrath u. s. w.) das Geldvermögen der Witwenkasse für den 1. Julius 1851 auf
\begin{center}121548 Thaler 22 Ggr.\end{center}
\noindent\hspace*{\parindent}Um, so weit ich dazu im Stande bin, die Reduction auf den 1. October 1851 abzuschätzen, ziehe ich von dieser Summe ab die auf Michaelis fällig gewordenen Witwenpensionen mit . 1787 Thl. 12 Ggr. und setze hinzu
\begin{center}
\small\setlength{\tabcolsep}{2pt}
\begin{tabular}{l r l r l}
\shortstack[l]{den Betrag der Beiträge für das Jahr 1850 bis 1851\\. . . . . . . . . . . . . . . . . . . . . .}&780&,,&&\\
\shortstack[l]{den halbjährigen Pachtzins für die Apotheke\\. . . . . . . . . . . . . . . . . . . . . . . . . . .}&500&,,&&\\
\shortstack[l]{und einen vierteljährigen Betrag der Zinseinnahme,\\unter Zugrundelegung der letztjährigen mit Abzug\\der Einziehungskosten . . . . . . . . . . . . . . . . . . . . . .}&1248&,,&10&,,
\end{tabular}
\end{center}
\noindent woraus dann obige Position (9) hervorgeht.

\begin{center}\rule{.28\linewidth}{.3pt}\end{center}

In Erwägung des bedeutenden Plus, mit welchem die Bilanzrechnung abschliesst, erscheint eine Erhöhung der Pension für die nächst bevorstehende Periode als zulässig, und über die Grösse der Erhöhung habe ich Folgendes zu bemerken.

Wenn die Frage aufgestellt wird, um wie viel unter Beibehaltung aller übrigen Einrichtungen der bewegliche Theil der Witwenpension erhöhet werden muss, damit in der Bilanz Debet und Credit zur vollkommenen Gleichheit gebracht werden, so findet sich durch eine leichte Rechnung diese Erhöhung \(=\) 13 Thl. 7 Ggr. Wählt man eine kleinere Erhöhung, so schliesst die veränderte Bilanz noch immer mit einem Plus ab, mit einem Minus hingegen, wenn eine grössere Erhöhung angenommen wird. Wurde also der bewegliche Theil der Pension von jetzt an auf 60 Thaler normirt, so dass jede der daran berechtigten Witwen zusammen 260 Thaler erhielte, so lange deren Anzahl nicht 18 überschreitet, im entgegengesetzten Falle hingegen neben dem festen Theil zu 200 Thaler noch den betreffenden Antheil an der Totalsumme 1080 Thaler, so würde die auf gleiche Art wie oben geführte Bilanzrechnung noch mit einem Plus von 1695 Thalern abschliessen. Man hat also zu dieser Maassregel nicht nur vollkommene Berechtigung, sondern auch die Aussicht, dass nach wenigen Jahren eine abermalige Erhöhung wird Statt finden können, insofern keine \textit{grosse} Verluste eintreten, der Genuss höhern Zinsfusses noch fortdauert, und bei der jetzt nicht erreichten Anzahl der Witwen 18 von der in Rechnung gebrachten jährlichen Summe vorerst jährlich etwas erübrigt wird.

Mit einem Minus von 877 Thalern hingegen würde die Bilanz abschliessen, wenn die Erhöhung auf 15 Thaler, und mit einem Minus von 3448 Thalern, wenn dieselbe auf 20 Thaler festgesetzt würde. Ein so geringes Minus, wie das im erstern Falle sich ergebende, würde aus den eben angeführten Gründen schon nach kurzer Zeit sich ausgleichen, und daher die Normirung des beweglichen Theils der Pension auf 65 Thaler (folglich bei mehr als 18 Witwen Vertheilung der Summe von 1170 Thalern) an sich gar kein Bedenken haben: vielleicht aber würde man nicht gern von dem bisher immer beobachteten Gebrauch abweichen wollen, wonach die halbjährige Pensionssumme stets eine ganze Anzahl Pistolen betragen hat (welche kleine Bequemlichkeit allerdings von selbst wegfallen wird, sobald die Anzahl der Witwen über 18 gestiegen ist). Schon jetzt aber den beweglichen Theil auf 70 Thaler zu setzen, würde ich schon des
